Una càrrega puntual $q_\mathrm{A} = -2,00\ \mathrm{\mu C}$ se situa a l'origen de coordenades, i una càrrega puntual $q_\mathrm{B} = -4,00\ \mathrm{\mu C}$ se situa a la posició $x = 1,00\ \mathrm{m}$ i $y = 0,00\ \mathrm{m}$.

\begin{apartat}{1,25}
Representeu les dues càrregues a la quadrícula (pla $xy$). Dibuixeu-hi també qualitativament (de manera aproximada) les línies del camp elèctric i les línies equipotencials generades per aquest sistema de dues càrregues.
\end{apartat}

\begin{apartat}{1,25}
En quin punt de l'eix $x$ s'ha de situar una tercera càrrega, $q_\mathrm{C} = -8,00\ \mathrm{\mu C}$, perquè la força elèctrica total sobre $q_\mathrm{C}$ sigui nul·la? Calculeu el treball que fa el camp elèctric quan la càrrega $q_\mathrm{C}$ es mou de la posició inicial que heu calculat fins a un punt equidistant de $q_\mathrm{A}$ i $q_\mathrm{B}$ situat a l'eix $x$. Interpreteu el signe del resultat obtingut.
\end{apartat}

\dades{%
  $k = \dfrac{1}{4\pi\varepsilon_0} = 8,99\times 10^{9}\ \mathrm{N\,m^2\,C^{-2}}$.}

\figura[0.55]{fig1.png}
